\section{Conclusion}
Having compute-intensive nature and heavy workload of data make the verification stage a bottleneck for the entire read mapping process. Pre-alignment filters are designed to facilitate verification by means of reducing the workload as accurately and fast as possible. Different techniques and hardware platforms are utilized for a quick filtering experience. In that sense, GateKeeper-GPU positions itself at a middle level of being the most accurate and fastest pre-alignment filtering tool. Compared to GateKeeper, which was implemented in FPGA, GateKeeper-GPU is more accurate but its benefit on the acceleration of the entire read mapping process is smaller. On the other hand, being a GPGPU tool makes it preferable to an FPGA tool. Since it is implemented on GPU, it is also much more promising for further improvements. 

We believe that, although GateKeeper-GPU still brings benefits, it is more advantageous to consider it when a new short read alignment tool is constructed with a verification-aware design rather than adapting it to an existing read mapping workflow since it requires extra steps for filtration such as encoding. With a well-developed hardware-software co-design of read mapping, it can have a powerful impact on the entire mapping procedure. For that reason, we provide GateKeeper-GPU with two different modes in encoding, both of which can be desirable in different scenarios. If the read mapper is designed to process encoded sequences in its original workflow, GateKeeper-GPU's encoding can be skipped, and the host-encoded version can be utilized. In other scenarios where the time spent for encoding can be overlapped with other operations and be hidden during kernel execution, a device-encoded version can be useful. 

The actual reasons that degrade GateKeeper-GPU's overall performance are tied to preprocessing work that has to be carried out for kernel execution. In future work, we intend to solve these problems and improve the extra time spent on the preprocessing to bring out GateKeeper-GPU's best performance. In addition, we notice that GateKeeper-GPU mainly utilizes L2 cache with an average hit rate of $86.2\%$ rather than unified/texture L1 cache. The hit rate of unified/texture L1 cache is $31.2\%$ on average, which is low. Improving cache utilization is another aim for enabling more efficient kernel activity. 

GateKeeper-GPU brings many benefits over its original work, even though the acceleration is less. It also has comparable results with similar tools. Having the issues resolved, the improvements can be more pronounced. Therefore, GateKeeper-GPU is a promising pre-alignment tool with a wide range of platform support owing to its simple design and GPGPU codebase.